\documentclass[11pt,a4paper]{article}

\usepackage[utf8]{inputenc}
\usepackage[T1]{fontenc}
\usepackage[french]{babel}
\usepackage{lmodern}
\usepackage{amsmath}
\usepackage{graphicx}
\usepackage{booktabs}
\usepackage{hyperref}
\usepackage{geometry}

\geometry{margin=2.5cm}

\title{Résolution du Problème du Voyageur de Commerce \\ par Métaheuristiques et Approche Hybride}
\author{Master Intelligence Artificielle et Recherche Opérationnelle \\ Loubna Elasri}
\date{Année universitaire : 2025--2026}

\begin{document}

\maketitle

\begin{abstract}

Le problème du voyageur de commerce (Traveling Salesman Problem, TSP) constitue l'un des problèmes les plus étudiés en optimisation combinatoire.
Il consiste à déterminer une tournée de coût minimal permettant de visiter un ensemble de villes exactement une fois avant de revenir au point de départ.

En raison de la complexité combinatoire du problème, les méthodes exactes deviennent rapidement inefficaces lorsque la taille des instances augmente.
Les métaheuristiques représentent donc une alternative efficace permettant d'obtenir des solutions de bonne qualité dans des temps de calcul raisonnables.

Dans ce projet, plusieurs approches sont étudiées pour résoudre le TSP : une heuristique gloutonne basée sur la méthode du plus proche voisin, l'algorithme Threshold Accepting et le recuit simulé.
Afin d'améliorer les performances obtenues, une approche hybride combinant le recuit simulé avec une recherche locale de type 2-opt est également proposée.

Les différentes méthodes sont testées sur plusieurs instances du problème et les résultats obtenus sont analysés et comparés.

\end{abstract}

\section{Introduction}

Le problème du voyageur de commerce (TSP) est un problème classique d'optimisation combinatoire.
Il apparaît dans de nombreux domaines tels que la logistique, la planification de tournées, la robotique ou encore la bio-informatique.

Le problème consiste à déterminer la plus courte tournée permettant de visiter un ensemble de villes une seule fois chacune, avant de revenir à la ville de départ.

Soit $n$ villes et une matrice de distances $D=(d_{ij})$ où $d_{ij}$ représente la distance entre la ville $i$ et la ville $j$.

Une solution du problème est une permutation des villes :

\[
x = (x_1,x_2,\dots,x_n)
\]

Le coût d'une tournée est défini par :

\[
f(x)=\sum_{k=1}^{n} d_{x_k x_{k+1}}, \quad \text{avec } x_{n+1}=x_1
\]

Le TSP appartient à la classe des problèmes NP-difficiles.
Cela signifie que le temps nécessaire pour trouver une solution optimale augmente de manière exponentielle avec le nombre de villes.

Pour cette raison, les heuristiques et les métaheuristiques sont largement utilisées afin d'obtenir des solutions proches de l'optimum.

Dans ce travail, plusieurs méthodes sont implémentées et comparées afin d'étudier leur efficacité pour la résolution du TSP.

\section{Lecture des données}

Les instances du problème sont fournies sous forme de matrices de distances stockées dans des fichiers texte.
Chaque élément $d_{ij}$ de la matrice représente la distance entre la ville $i$ et la ville $j$.

Lors de la lecture des données, la matrice est stockée dans une structure de type tableau permettant un accès rapide aux distances.
Une vérification est également effectuée afin de garantir que $d_{ii}=0$ pour toute ville $i$.

\section{Objectif du projet}

Le sujet demande d'implémenter :

\begin{itemize}
\item une heuristique de base du \emph{plus proche voisin} ;
\item une métaheuristique pour améliorer la solution initiale ;
\item une comparaison des performances en termes de coût de tournée et de temps d'exécution.
\end{itemize}

\section{Heuristique gloutonne du plus proche voisin}

La première méthode utilisée est l'heuristique gloutonne du plus proche voisin.

\begin{enumerate}
\item Choisir une ville de départ.
\item Sélectionner la ville non visitée la plus proche.
\item Ajouter cette ville à la tournée.
\item Répéter jusqu'à ce que toutes les villes soient visitées.
\item Revenir à la ville de départ.
\end{enumerate}

Cette heuristique est très rapide mais peut conduire à des solutions de qualité limitée.

\section{Algorithme Threshold Accepting}

Threshold Accepting est une méthode de recherche locale.

À chaque itération un voisin est généré en échangeant deux villes dans la tournée.
Si la variation de coût est inférieure à un seuil $S$, la solution est acceptée.

Le seuil est progressivement diminué au cours de l'exécution.

\section{Recuit simulé}

Le recuit simulé est inspiré du processus physique de refroidissement.

Une solution moins bonne peut être acceptée avec une probabilité :

\[
P = e^{-d/T}
\]

où $d$ représente la détérioration de la solution et $T$ la température.

La température diminue progressivement :

\[
T \leftarrow \alpha T
\]

avec $0<\alpha<1$.

\section{Approche hybride}

L'approche hybride combine le recuit simulé avec la recherche locale 2-opt.

La méthode 2-opt consiste à supprimer deux arêtes de la tournée et reconnecter les segments afin de réduire la distance totale.

\section{Résultats expérimentaux}

\begin{table}[h]
\centering
\begin{tabular}{lcccc}
\toprule
 & Greedy & Threshold Accepting & SA & Hybrid \\
\midrule
Distance & 2394 & 2220 & 2157 & 2050 \\
Temps (s) & 0.0 & 2.257 & 2.498 & 4.415 \\
\bottomrule
\end{tabular}
\caption{Comparaison des méthodes pour l'instance 29}
\end{table}

\begin{table}[h]
\centering
\begin{tabular}{lcccc}
\toprule
 & Greedy & Threshold Accepting & SA & Hybrid \\
\midrule
Distance & 558.85 & 479.81 & 479.81 & 490.55 \\
Temps (s) & 0.0 & 2.428 & 3.385 & 5.713 \\
\bottomrule
\end{tabular}
\caption{Comparaison des méthodes pour l'instance 51}
\end{table}

\section{Conclusion}

Les résultats montrent que les métaheuristiques améliorent la solution gloutonne initiale.

L'approche hybride combinant recuit simulé et recherche locale 2-opt permet d'obtenir les meilleures performances.

\end{document}
